\subsection{\texorpdfstring{$S_4$}{S4}-Hurwitz Jacobian 的最小同源基、隐含商曲线阴影与 Boolean 下闭压缩的 dyadic 刚性}\label{subsec:conclusion-s4-hurwitz-basis-boolean-dyadic-rigidity}
沿用
\[
H:=X/A_4,\qquad
Z:=X/D_8,\qquad
CF:=X/S_3,\qquad
Y:=X/C_4,\qquad
P:=\operatorname{Prym}(Y/Z)
\]
的记号，并记
\[
\mathcal B:=\{J(H),\,J(Z),\,J(CF),\,P\}.
\]
此前结论已经分别给出 \(S_4\)-Hurwitz 商塔 Jacobian 的四生成元完备性与极小性、Chevalley--Weil 重数与 Prym 因子的纯化对应、以及 Boolean 互斥矩阵在任意下闭主子矩阵上的精确行列式与 Smith 形。将这些接口并置后，还可继续抽出下列两条彼此正交的终端刚性：一条是由 \(\mathcal B\) 控制的最小同源基结构，另一条是 Boolean 下闭压缩所满足的严格 dyadic 剩余律。

\begin{theorem}[\texorpdfstring{$S_4$}{S4}-Hurwitz Jacobian 的最小同源基与表示纯化]\label{thm:conclusion-s4-hurwitz-minimal-isogeny-basis-purification}
在由全部 subgroup-quotient Jacobians
\[
\{J(X/U):\ U\le S_4\}
\]
生成的 \(\QQ\)-同源范畴中，四元组 \(\mathcal B\) 构成极小生成基。更具体地，
\[
J(X)\sim_{\QQ}J(H)\times J(Z)^2\times J(CF)^3\times P^3,
\]
并且 \(J(X)\) 的四个非平凡 \(S_4\)-等型部分满足
\[
A_{\mathrm{sgn}}\sim_{\QQ}J(H),\qquad
A_{V_2}\sim_{\QQ}J(Z)^2,
\]
\[
A_{V_3}\sim_{\QQ}J(CF)^3,\qquad
A_{V_3'}\sim_{\QQ}P^3.
\]
因此，\(\mathcal B\) 的四个因子恰好逐一纯化
\[
5\cdot \mathrm{sgn},\qquad
4\cdot V_2,\qquad
3\cdot V_3,\qquad
9\cdot V_3'
\]
这四个非平凡 Hodge 块；特别地，任何真子集 \(\mathcal B'\subsetneq \mathcal B\) 都不能生成全部 subgroup-quotient Jacobians 的有理同源类型。
\end{theorem}
\begin{proof}
四生成元完备性与极小性已由定理 \ref{thm:conclusion-s4-hurwitz-jacobian-four-generator-semiring} 与定理 \ref{thm:conclusion-s4-hurwitz-jacobian-generator-minimality} 给出。另一方面，定理 \ref{thm:xi-time-part9n1b-prym-chevalley-weil-reciprocity} 把
\[
H^0(X,\Omega_X^1)\cong
5\cdot \mathrm{sgn}\oplus 4\cdot V_2\oplus 3\cdot V_3\oplus 9\cdot V_3'
\]
与
\[
J(H),\qquad J(Z),\qquad J(CF),\qquad P
\]
之间的等型对应完全纯化为
\[
A_{\mathrm{sgn}}\sim_{\QQ}J(H),\qquad
A_{V_2}\sim_{\QQ}J(Z)^2,\qquad
A_{V_3}\sim_{\QQ}J(CF)^3,\qquad
A_{V_3'}\sim_{\QQ}P^3.
\]
于是四个几何因子既构成全部商 Jacobian 的极小同源基，又逐一隔离四个非平凡不可约块，因而彼此不可替代。证毕。
\end{proof}

\begin{corollary}[隐含商曲线亏格的仿射阴影]\label{cor:conclusion-s4-hidden-quotient-genera-affine-shadows}
\leanverified{paper\_conclusion\_s4\_hidden\_quotient\_genera\_affine\_shadows}
一切 subgroup-quotient genera 都是
\[
\bigl(g(H),\,g(Z),\,g(CF),\,\dim P\bigr)
\]
的整系数仿射阴影。若记
\[
C_6:=X/V_4,
\]
则特别地，由
\[
J(Y)\sim_{\QQ}J(Z)\times P,\qquad
J(C_6)\sim_{\QQ}J(H)\times J(Z)^2
\]
可得
\[
g(Y)=g(Z)+\dim P,\qquad
g(C_6)=g(H)+2g(Z),
\]
从而两个隐含商曲线亏格满足
\[
g(X/C_2)=2g(CF)+g(Y)=2g(CF)+g(Z)+\dim P=19,
\]
\[
g(X/C_3)=\frac{3g(H)+2g(CF)+2g(Y)-g(C_6)}{2}
=g(H)+g(CF)+\dim P=17.
\]
因此，\(g(X/C_2)\) 与 \(g(X/C_3)\) 不是新的独立不变量，而只是最小同源基维数数据的仿射投影。
\end{corollary}
\begin{proof}
由定理 \ref{thm:conclusion-s4-hurwitz-jacobian-four-generator-semiring}，
\[
J(X/C_2)\sim_{\QQ}J(Z)\times J(CF)^2\times P,
\qquad
J(X/C_3)\sim_{\QQ}J(H)\times J(CF)\times P,
\]
故直接取维数便得到
\[
g(X/C_2)=g(Z)+2g(CF)+\dim P,
\qquad
g(X/C_3)=g(H)+g(CF)+\dim P.
\]
另一方面，定理 \ref{thm:xi-time-part9n1b-kani-rosen-hidden-quotient-genera} 也给出了同一对象的 Burnside--Kani--Rosen 阴影公式
\[
g(X/C_2)=2g(CF)+g(Y),
\]
\[
g(X/C_3)=\frac{3g(H)+2g(CF)+2g(Y)-g(C_6)}{2}.
\]
再由
\[
J(Y)\sim_{\QQ}J(Z)\times P,\qquad
J(C_6)\sim_{\QQ}J(H)\times J(Z)^2
\]
取维数，即得
\[
g(Y)=g(Z)+\dim P,\qquad
g(C_6)=g(H)+2g(Z).
\]
代入
\[
g(H)=5,\qquad g(Z)=4,\qquad g(CF)=3,\qquad \dim P=9
\]
便得到
\[
g(X/C_2)=4+2\cdot 3+9=19,\qquad
g(X/C_3)=5+3+9=17.
\]
证毕。
\end{proof}

\begin{theorem}[Boolean 下闭压缩的 dyadic unimodularity]\label{thm:conclusion-boolean-orderideal-dyadic-unimodularity}
\leanverified{paper\_conclusion\_boolean\_orderideal\_dyadic\_unimodularity}
设
\[
V_q:=2^{[q]},
\qquad
A_q:=B^{(q)}(2,1).
\]
对任意下闭集 \(I\subseteq V_q\)，记 \(A_q[I]\) 为按 \(I\) 取行列得到的主子矩阵。则总有
\[
\det A_q[I]\in\{\pm 1,\pm 2\}.
\]
因此 \(A_q[I]\) 始终是非奇异整数矩阵，并且
\[
A_q[I]^{-1}\in \frac12 M_{|I|}(\ZZ).
\]
换言之，全部下闭压缩共同构成一个严格的 dyadic unimodular 家族。
\end{theorem}
\begin{proof}
推论 \ref{cor:xi-disjointness-boolean-order-ideal-principal-minors} 已给出
\[
\det A_q[I]
=2^{\ind_{\{\varnothing\in I\}}}
(-1)^{\#\{S\in I\setminus\{\varnothing\}:\ \card{S}\ \mathrm{odd}\}},
\]
故 \(|\det A_q[I]|\in\{1,2\}\)。于是 \(A_q[I]\) 必为非奇异整数矩阵。再由伴随矩阵公式
\[
A_q[I]^{-1}=\frac{\operatorname{adj}(A_q[I])}{\det A_q[I]},
\]
便立得
\[
A_q[I]^{-1}\in \frac12 M_{|I|}(\ZZ).
\]
证毕。
\end{proof}

\begin{theorem}[Boolean 下闭压缩的单比特奇偶缺陷]\label{thm:conclusion-boolean-orderideal-onebit-parity-defect}
\leanverified{paper\_conclusion\_boolean\_orderideal\_onebit\_parity\_defect}
沿用定理 \ref{thm:conclusion-boolean-orderideal-dyadic-unimodularity} 的记号。则对任意非空下闭集 \(I\subseteq V_q\)，\(A_q[I]\) 的 Smith 形恰落在下列三类之一：
\[
\mathrm{SNF}(A_q[I])=
\operatorname{diag}(1,\dots,1),
\]
\[
\mathrm{SNF}(A_q[I])=\operatorname{diag}(2),
\]
或
\[
\mathrm{SNF}(A_q[I])=\operatorname{diag}(1,\dots,1,2).
\]
等价地，
\[
\operatorname{coker}(A_q[I])
\cong 0
\qquad\text{或}\qquad
\operatorname{coker}(A_q[I])\cong \ZZ/2\ZZ.
\]
因此，任意整数方程
\[
A_q[I]x=y,\qquad y\in\ZZ^{|I|}
\]
的整解存在性至多受一个单独的 parity bit 阻断；既不会出现更高 \(2\)-幂挠性，也不会出现任何奇素数挠性。
\end{theorem}
\begin{proof}
在定理 \ref{thm:xi-boolean-two-layer-order-ideal-smith-shape-invariance} 中取
\[
a=2,\qquad b=1,\qquad d=a-b=1,\qquad g=\gcd(a,b)=1,
\]
便得到：
若 \(\varnothing\notin I\)，则
\[
\mathrm{SNF}(A_q[I])=\operatorname{diag}(1,\dots,1);
\]
若 \(I=\{\varnothing\}\)，则
\[
\mathrm{SNF}(A_q[I])=\operatorname{diag}(2);
\]
若 \(\varnothing\in I\) 且 \(|I|\ge 2\)，则
\[
\mathrm{SNF}(A_q[I])=\operatorname{diag}(1,\dots,1,2).
\]
于是余核只能是平凡群或 \(\ZZ/2\ZZ\)。这说明全部 order-ideal 压缩的整数剩余仅保留一个 dyadic 奇偶位，而不存在更高阶或奇素数方向。证毕。
\end{proof}

\begin{corollary}[Boolean 互斥 calculus 的二元 residual-axis]\label{cor:conclusion-boolean-orderideal-residual-axis-dichotomy}
\leanverified{paper\_conclusion\_boolean\_orderideal\_residual\_axis\_dichotomy}
将 \(\operatorname{coker}(A_q[I])\) 视为有限 residue 对象，则其连续圆维--pro 圆维二元签名只能取
\[
\Sigma_{\circ,\mathrm{pro}}\bigl(\operatorname{coker}(A_q[I])\bigr)
=(0,0)
\qquad\text{或}\qquad
\Sigma_{\circ,\mathrm{pro}}\bigl(\operatorname{coker}(A_q[I])\bigr)=(0,1).
\]
因此，对任意下闭压缩的精确外置，至多需要一个 dyadic residue 轴，而从不需要奇素数 register。
\end{corollary}
\begin{proof}
由定理 \ref{thm:conclusion-boolean-orderideal-onebit-parity-defect}，
\[
\operatorname{coker}(A_q[I])\cong 0
\qquad\text{或}\qquad
\operatorname{coker}(A_q[I])\cong \ZZ/2\ZZ.
\]
而命题 \ref{prop:cdim-finite-fiber-min-prime-register-axes} 与备注 \ref{rem:cdim-cont-pro-signature} 表明：对有限交换群 \(F\)，
\[
\Sigma_{\circ,\mathrm{pro}}(F)=(0,d(F)).
\]
其中 \(d(0)=0\)，\(d(\ZZ/2\ZZ)=1\)。故只可能得到
\[
(0,0)\qquad\text{或}\qquad(0,1).
\]
证毕。
\end{proof}

\begin{conjecture}[\texorpdfstring{$S_4$}{S4}-Hurwitz 最小同源基与 Boolean dyadic 缺陷的桥接]\label{conj:conclusion-s4-boolean-prym-dyadic-bridge}
设
\[
\mathcal I_q:=\{I\subseteq V_q:\ I\ \text{下闭}\},
\]
并记
\[
\langle \mathcal B\rangle_{\QQ\text{-isog}}
\]
为由 \(\mathcal B=\{J(H),J(Z),J(CF),P\}\) 生成的 \(\QQ\)-同源张成。则应存在一个 functorial bridge
\[
\mathfrak F_q:\mathcal I_q\longrightarrow \langle \mathcal B\rangle_{\QQ\text{-isog}}
\]
满足：
\begin{enumerate}
  \item \(\operatorname{coker}(A_q[I])\cong 0\) 当且仅当 \(\mathfrak F_q(I)\) 不含 \(V_3'\)-方向；
  \item \(\operatorname{coker}(A_q[I])\cong \ZZ/2\ZZ\) 当且仅当 \(\mathfrak F_q(I)\) 只在 \(P\)-方向上保留最小非平凡剩余；
  \item 因而 Boolean 互斥 calculus 的单比特奇偶缺陷，应当正是 \(S_4\)-Hurwitz 链中 \(V_3'\)-Prym 块的离散 shadow。
\end{enumerate}
\end{conjecture}

\endinput
